\scanpage{218}

\mspage{92}

\[
\widetilde C=\widetilde C_0\tag*{(7.3.3)}
\]

(Ce lemme implique aussitôt la première \unsure{assertion} de 7.1. $1^\circ$.) Prouvons ce lemme, \deleted{\unsure{ou}} \inserted{i.e. que \deleted{tout} $i\in C$ est dans $\widetilde C_0$,} en construisant par récurrence transfinie \inserted{une suite} \deleted{des} \inserted{\unsure{str.|str} croissante de} sous-objets de $B$,

\[
A_\alpha\overset{j_\alpha}{\lhook\joinrel\relbar} B,\qquad \text{avec \inserted{les} }j_\alpha\in C,\ \text{et}
\]

\inserted{avec} \deleted{de telle façon que} \textcircled{$1^\circ$} $A_0=A$, \deleted{et} \inserted{On prouvera que} $A\xrightarrow{i_\alpha=i_{\alpha,0}}A_\alpha$ dans $\widetilde C_0$, \deleted{\unsure{et de}} et \deleted{Si $j$ est un iso} il y aura un ordinal final $\alpha$ tel que $A_\alpha=B$; \deleted{\unsure{et}} on en conclura donc que $i=i_\alpha$ est dans $\widetilde C_0$.

Si $i$ est un iso, on sait que $i\in\widetilde C_0$, il \deleted{\unsure{On doit}} n'y a rien à prouver (la constr. transfinie s'arrête au cran 0). Les \deleted{deux} \inserted{deux} pas de la construction, après le pas initial $1^\circ$, sont :

\textcircled{$2^\circ$} Soit $A_\alpha$ construit. Si $A_\alpha=B$, la construction est arrêtée. Si $A_\alpha\xrightarrow{j_\alpha} B$ \unsure{non iso|un iso}, on applique 7.3.1. à $j_\alpha\in C$, pour trouver un \deleted{\unsure{$-$}} $C\overset{j}{\hookrightarrow}B$, avec $j\in C$. et $C\supsetneqq A_\alpha$. De plus,

\[
i_{\alpha+1,\alpha}\colon A_\alpha\mathrel{\unsure{-|\to}} A_{\alpha+1}\overset{\mathrm{def}}=C\ \text{est dans }\widetilde C_0,\tag*{(7.3.4.)}
\]

car déduit de $i'\in\widetilde C_0$ par cochgt. de base.

\textcircled{$3^\circ$} Si $\alpha$ est ordinal limite, on pose

\[
A_\alpha=\varinjlim_{\beta<\alpha} A_\beta\tag*{(7.3.5)}
\]
% archive sheet number: 217
